\input preamble.tex
\begin{centering}
\Huge{\textbf{Prøve 01 i strømningsmåling}}\\
\end{centering}
\vskip 2cm 
Følgende kompetansemål er relevante for prøven:
\begin{itemize}[noitemsep]

	\item planlegge, utføre, vurdere kvalitet, sluttkontrollere og dokumentere arbeidet
	\item montere, konfigurere, kalibrere og idriftsettelse digitale og analoge målesystemer
\end{itemize}

Alle ark som leveres inn skal ha elevens navn.
\vskip 2.5pt 
Oppgave 1-8 skal leveres på papir, etter levering kan eleven ta frem PC og svare på oppgave 9.\\
\vskip 2.5pt 
Når oppgave 9 skal leveres kan elevene slå på trådløst nettverkt og sende oppgaven på mail til fred-olav.mosdal@skole.rogfk.no
\vskip 2.5pt 
Når de bruker PC er alle ikke kommuniserende hjelpemidler tillatt
\vskip 2.5pt
Konaktinformasjon:
\begin{itemize}[noitemsep]
	\item Kontaktlærer: Fred-Olav Mosdal
	\item TLF: 90507684
\end{itemize}


\vfil \eject
%Slutt forside
Oppgave 1 (6p)%Navngi
\vskip 2.5pt 
a) Nevn 4 hastighetbaserte strømningsmålere. \\
\vskip 2.5pt 
\begin{tikzpicture}
	\draw[step=0.5cm,gray!20,very thin]  grid (17,4) ;
\end{tikzpicture}
\vskip 2.5pt 
b) Nevn 4 strømningsmålere som er basert på differansetrykk\\
\vskip 2.5pt 
\begin{tikzpicture}
	\draw[step=0.5cm,gray!20,very thin]  grid (17,4) ;
\end{tikzpicture}
\vskip 2.5pt 


Oppgave 2 (6p) %Definisjoner
\vskip 2.5pt 
a) Hva forteller Reynolds nummer oss? \\
\vskip 2.5pt 
\begin{tikzpicture}
	\draw[step=0.5cm,gray!20,very thin]  grid (17,6) ;
\end{tikzpicture}
\vskip 2.5pt 

b) Hva menes med massestrøm sett i forhold til volumstrøm? \\
\vskip 2.5pt 
\begin{tikzpicture}
	\draw[step=0.5cm,gray!20,very thin]  grid (17,6) ;
\end{tikzpicture}
\vskip 2.5pt 
\vskip 2.5pt 
\vfil\eject
Oppgave 3 (6p)%Enkel utregninger.
\vskip 2.5pt 
a) Convert the volumetric flow rate of 35 gallons per minute (35 GPM) into a {\it mass} flow rate in pounds per minute, assuming the fluid in question is water.\\

\vskip 2.5pt 
\begin{tikzpicture}
	\draw[step=0.5cm,gray!20,very thin]  grid (17,6) ;
\end{tikzpicture}
\vskip 2.5pt 
b) Vi har et rør som det strømmer olje med en strømningsrate på 200 m³/h og en temperatur på 50°C. Begge seksjonene er etter schedule 40.  Den første delen av røret har dimensjon DN200 (ID=202.74mm) og den andre delen har dimensjon DN65 (ID=62.68mm)
$$\includegraphics[width=15.5cm]{aFlow01x02.eps}$$
Regn ut hastigheten til fluidet i hver av seksjonene. 
\vskip 2.5pt 
\begin{tikzpicture}
	\draw[step=0.5cm,gray!20,very thin]  grid (17,12) ;
\end{tikzpicture}
\vskip 2.5pt 
\vskip 2.5pt 
Oppgave 4 (6p) %Tegneoppgaver(blokkskjema skisse osv. 
\vskip 2.5pt 
a) Tegn en laminær strømningsprofil \\
\vskip 5cm
b) Tegn en turbulent strømningsprofil \\
\vskip 5cm 
% Her kommer oppgaver som skal kreve forklaringer eller forståelse
\vfil\eject
Oppgave 5 (6p) %Feilsøkings oppgave
\vskip 2.5pt 
What will be wrong with this measurement system if we connect a linear-scale indicator (an electrical meter movement responding to the transmitter's current signal) to the transmitter's output, and try to measure fluid flow along this scale?  Assume the transmitter has been properly calibrated to output full current (typically 20 mA) at full flow through the orifice plate.\\
$$\includegraphics[width=10.5cm]{aFlow01x01.eps}$$
\\
\begin{tikzpicture}
	\draw[step=0.5cm,gray!20,very thin]  grid (17,9) ;
\end{tikzpicture}
\vskip 2.5pt 
\vfil \eject
Oppgave 6 (6p) % Tegn og forklar virkemåte
\vskip 2.5pt 
Tegn og forklar virkemåten til en måleblende. 
\vfil \eject

Oppgave 7 (6p) % Tegn og forkalr virkemåte. 
\vskip 2.5pt 
Tegn og forklar virkemåten til en magnetiske strømningsmåler. 
\vfil \eject
Oppgave 8 (6p) % Regneoppgave
\vskip 2.5pt 
Vi har en vortex strømnigsmåler som måler måler drivstofforbruket til en  forbrenningsovn. Strømningsmåleren har en "k=2.33liter per puls"
\vskip 2.5pt 

The sensor frequency at a fuel oil flow rate of 8510 liter per hour.

\vskip 10pt

The total amount of fuel consumed by the boiler after a digital counter circuit records 800,000 pulses.

\vskip 10pt

The fuel oil flow rate (in cubic feet per minute) at a sensor frequency of 35 Hz.

\vskip 10pt
\vfil\eject
Oppgave 9 (12p) \\ %skal beskrive hvordan en jobb skal utføres. (Planlegg, besriv hvordan du ville gjenomført og dokumenter jobben)
Du jobber i et firma som blant annet driver med produksjon av aceton. I forbindelse med destilasjonskollonen har det vært problemer med strømningmåleren på reflux linjen. Den er av type Rousemount Model 8800C Smart Vortex Flowmeter. \\
Temperaturen til acetonet i refluxlinjen er 100°C. Ved denne temperaturen har aceton en $\rho=693.7 kg/m³$ og en µ=0.17mPa s\\
$$\includegraphics[width=15.5cm]{i04060x01.eps}$$
Måleren tatt med inn til benksetting (bench commissioning, se side 2-2 i manual. Du skal utføre det som er beskrevet under "A configure". Jobben din er ferdig når du er kommet til B). Din leder ønsker at du setter URV så høyt som måleren kan måle med Aceton under de gitte betingelsene, og at du setter LRV så lavt som det er mulig. \\
\\
Rørlinjen som måleren står på er et DN25 Schedule 40 rør med ID=26.64mm\\\\
Beskriv hvordan du ville planlagt, gjennomført og dokumentert denne jobben.

\end {document}
